

\section{Assessment}
Since this project is a group tasks, each group is awarded a Collective Group Mark on the standard university scale, with the different aspects of the project contributing according to the following table:


\begin{table}[h!]
    \centering
    \footnotesize
    \renewcommand{\arraystretch}{1.2} % Adjust row height for readability
    \caption{Assessment Details}
    \begin{tabular}{p{4cm} p{2cm} p{8cm}} % Define column widths
    \hline
    \textbf{Assessment Type} & \textbf{Weight} & \textbf{Detail} \\ \hline
    Practical Test & 30\% & This is the individual report that shall be submitted every week. \\
    Multimedia Recording (Report) & 20\% & The group will try to communicate and sell/pitch their concept and will be evaluated based on their communication and pitching style. \\
    Report Writing (Assignment) & 40\% & Each group will write a report and summarize what they had learned and improved, collectively. More detail about the report writing is available elsewhere. \\
    Presentation (In class) & 10\% & Here, we will evaluate the individual communication skill (whether they can give a proper idea/contribution) and whether they can give ideas that reflect their entrepreneurial skill. \newline This evaluation will be conducted randomly at any lab session. \newline \textbf{Marks are given individually.} \\
    Mini Project & 0\% & This is our grand event, where the group will pitch their invention in front of an audience and will be bombarded with questions (hopefully not). \newline \textbf{Students, as a group, need to work together to convince the "dragon" (hypothetical potential investor).} \\
    \hline
    \end{tabular}
    \label{tab:assessment}
\end{table}



\subsection{Peer Assessment}
Peer assessment is then used to derive an Individual Mark for Group Work, on the university scale, for each group member from the Collective Group Mark. This is done by distributing the collective mark among the group members according to merit as perceived by the group members themselves (but maybe vetted by the supervisor), such that the average of the individual marks for the group work equals the assigned Collective Group Mark (but maybe nomalized by the module convenor).

As explained above, peer assessment is a key aspect of the overall assessment. It is used to distribute the collective group mark to each individual according to merit as perceived by the peers, i.e. the other group members. Each group member is asked to rate each of his or her peers according to a number of aspects on a purpose-designed form. A written justification supporting the ratings is also required for each member. The peer assessments are submitted in complete confidence as part of the Individual Report, and will not be seen by the other students. Note that, under certain circumstances, it is possible to revise the assessment together with the module convenor. See below.


\subsection{How to access the Peer Assessment form}
The Peer Assessment Form can be found on the module's main page.

Be sure to read the instructions on the forms carefully, as well as the instructions below, before filling out the forms. The completed forms are to be included as an appendix in your Individual Report.
Each group member is assessed according to the following aspects:
\begin{itemize}
    \item Research and information gathering
    \item Creative input
    \item Co-operation within group
    \item Communication within group
    \item Quality and quantity of concrete contributions to group deliverables
    \item Attendance at meetings
\end{itemize}

Thus it is not just the amount of contributed work that is being evaluated, say in terms of words in the reports, but a wide spectrum of contributions all of which are important for a successful end result. This also means that everyone gets a chance to contribute according to their specific skills and get proper recognition for this. Clearly, someone who excels in designing the system architecture and implementing it has made a very important contribution. Likewise, someone who is a great writer and took the lead on getting the group reports together has also made a big contribution. But someone who was instrumental in mediating between conflicting views, say, and thus helped holding the group together has clearly also made a very valuable contribution to the group as a whole.
Each aspect is rated on a five-point scale ranging from “None” via “Adequate” to “Excellent”, where “Adequate” means that the assessed person has done what is expected for that aspect: no less, no more. Note that it is the relative performance of group members that matters as the peer marking only serves to redistribute the assigned collective group mark. For example, if everyone rates everyone else as “Excellent”, the end result is that everyone gets the same mark for their contribution to the group work which is going to be equal to the collective group mark.


Assessing your peers is a privilege and a big responsibility. Be fair and objective in your evaluations. In my previous experience, the member of my faculty have been very impressed with the quality and honesty of the peer assessment, and there has almost never been any problems.

\textbf{Important!} You need to provide a justification about the rating your provide. This is especially important if some of you try to play the system!!!!!
It worth to mention that, the module convener will carrying out a sanity check on the peer marks based on what they have learnt during the interaction with the group members throughout the 14 weeks, as well as from the group reports. Should there be obvious problems with some of the peer marks, then the module convener will override any or all of those marks where necessary.


\subsection{Revised Peer Assessment}
Note that the initial peer assessment is submitted as part of the individual reports. While this is sometime before the Presentation days, the majority of work will have been completed by then, and the time of report submission is thus a reasonable time to carry out the peer assessment.
However, it is conceivable that the performance of a peer during the final preparations for and at the Presentation days is significantly better or worse than his or her previous performance, so much so that the initial assessment becomes very unfair. Here “very unfair” means that the overall change of the assessment must be at least 3 points (three different aspects changed one step up, one aspect changed two steps down and one other aspect changed one step down, etc.). Should this be the case, it is permissible to revise the assessment for a peer one time, up until the end of the Presentation Day. The revised assessment is subject to vetting by the module convenor just like the initial one.
To revise an assessment, the following procedure must be followed:
\begin{itemize}
    \item The original assessment form for the peer in question should be amended to show the new assessment. It should also state that this is a revised assessment. Further, each individual change must be justified in writing by adding a paragraph that states the old assessment, the new assessment, and why it has changed.
    \item The assessment must be conducted and a copy of the form must be e-mailed in PDF format to the module convener no later than 18:00 on the Presentation Day.
\end{itemize}


